\documentclass[10pt,a4paper]{amsart}
\usepackage[margin=19mm]{geometry}
\usepackage{amsmath,amssymb,mathtools}
\usepackage[scr=rsfs,cal=euler]{mathalfa}
\usepackage{xfrac}
\usepackage[unicode,hidelinks,bookmarks=false]{hyperref}
\newcommand{\ie}{i.e.\ }
\newcommand{\cf}{cf.\ }
\newcommand{\resp}{respectively\ }
\newcommand{\Subsection}{\subsection}
\providecommand{\nobreakdash}{\nobreak\hbox{-}}
\setlength{\emergencystretch}{2em}
\multlinegap0pt
\usepackage{bm}
\usepackage{amssymb}
\usepackage{float}
\usepackage{dsfont}
\usepackage{mathtools}
\usepackage{xcolor}
\usepackage[shortlabels]{enumitem}
\newtheorem{assumption}{Assumption}
\newtheorem{theorem}{Theorem}
\newcommand{\bpi}{\boldsymbol{\pi}}
\newcommand \la {\lambda}
\title{Equilibrium Strategies for the N-agent Mean-Variance Investment Problem over a Random Horizon}
\author{Xiaoqing Liang, Jie Xiong, Ying Yang}
\date{Source: arXiv:2507.04611v2 (CC BY 4.0)}
\begin{document}
\maketitle

\noindent\emph{Source and licence.} Equilibrium Strategies for the N-agent Mean-Variance Investment Problem over a Random Horizon. Version arXiv:2507.04611v2 (CC BY 4.0), distributed by arXiv under the Creative Commons Attribution 4.0 International licence, http://creativecommons.org/licenses/by/4.0/, as recorded in the arXiv licence metadata for this exact version and shown on the abstract page https://arxiv.org/abs/2507.04611v2. Full licence notice: \texttt{licenses/arxiv-2507.04611-CC-BY-4.0.txt}.\par\smallskip

\noindent\textit{Editorial scope.} Counted unit: the article's theorem eq:M-pi\_i (source0.tex lines 955--980). The article's complete original proof of it is delivered with it (source0.tex lines 1481--1529, one \texttt{proof} environment of 49 lines, which the article places away from the statement). No mathematics is written, repaired or re-derived: the statement and the proof are the article's own lines, in the article's own notation, and no retained line is substituted or re-flowed.  Retained uncounted as the source-local context the proof needs: lines 923--927; lines 929--932; lines 934--938; lines 945--948; lines 950--953; lines 1024--1026.  Nothing was omitted from any retained span, so the package carries no omission record.  No retained line contains a citation, so the wrapper carries no bibliography and no bibliography entry.  No retained line contains a figure environment, an image-inclusion command or drawing code, so no image is delivered and the package declares no asset dependency.  Editorial formatting only, in addition: an AMS \texttt{amsart} preamble that restates the article's own declarations and macros used by the retained lines, namely the environments \texttt{assumption}, \texttt{theorem} on separate counters, together with the standard packages those lines need.  The retained environments print in this document as Assumption 1, Theorem 1 in order of appearance, whereas the article prints the same items with the numbering of its own full document.  The complete native source of the article as distributed in the arXiv e-print for this exact version is bundled unchanged as \texttt{sources/181108/source0.tex}.
	\begin{align}
		a_{\scriptscriptstyle m}=&\frac{\lambda}{\lambda-r-\rho_{\scriptscriptstyle m} p_{\scriptscriptstyle m}(b-r)},\label{eq:MFa1}\\
		 \tilde{a}_{\scriptscriptstyle m}=&\frac{\lambda }{\lambda-[2r+2\rho_{\scriptscriptstyle m} p_{\scriptscriptstyle m}(b-r)+\frac{1}{2}\rho_{\scriptscriptstyle m} {p_{\scriptscriptstyle m}}^2]}.\label{eq:MFbara}\\
		\rho_{\scriptscriptstyle m}=&\frac{1}{2(\xi^2+\sigma^2)}, \label{eq:MFrho}
	\end{align}

	\begin{align}	
		&\big(\gamma-\mu_{1}/2\big){z^3_{\scriptscriptstyle m}}+\big\{\mu_{1}\big[ (\lambda-r)/2-2\rho_{\scriptscriptstyle m}(b-r)^2\big]+\gamma[2\rho_{\scriptscriptstyle m}(b-r)^2-\lambda+2r]\big\}{z^2_{\scriptscriptstyle m}}\\
		&+\{\rho_{\scriptscriptstyle m}(b-r)^2[\mu_{1}(3\lambda-4r)+4r\gamma]+\gamma(\lambda-r)^2\}z_{\scriptscriptstyle m}+\rho_{\scriptscriptstyle m}(b-r)^2[2\gamma r^2-\mu_{1}(\lambda-r)(\lambda-2r)]=0.\\\label{eq:MF-p1}
	\end{align}	

	\begin{align}
		k_{1,\scriptscriptstyle m}&:={Q^{-1}_{\scriptscriptstyle m}}\rho_{\scriptscriptstyle m}\sigma(r-\lambda)(D_{\scriptscriptstyle m}-\mu_{1}c_{\scriptscriptstyle m}-2\gamma a_{\scriptscriptstyle m}c_{\scriptscriptstyle m}),\label{eq:MFk1}\\	
		k_{2,\scriptscriptstyle m}&:={Q^{-1}_{\scriptscriptstyle m}}\rho_{\scriptscriptstyle m} (b -r)[2\gamma (a_{\scriptscriptstyle m} -1)c_{\scriptscriptstyle m} +\mu_{1}c_{\scriptscriptstyle m} -D_{\scriptscriptstyle m} ],\label{eq:MFk2}\\
		k_{3,\scriptscriptstyle m}&:=-{Q^{-1}_{\scriptscriptstyle m}}\mu_{2}\rho_{\scriptscriptstyle m} \lambda(b -r).\label{eq:MFk3}
	\end{align}

	\begin{align}
		&\{\gamma \tilde{a}_{\scriptscriptstyle m}(\lambda-r)[2r-\lambda+(b-r)\rho_{\scriptscriptstyle m} p_{\scriptscriptstyle m}]-[2\gamma\lambda(a_{\scriptscriptstyle m}-1)+\mu_1(\lambda-r)]a_{\scriptscriptstyle m}\rho_{\scriptscriptstyle m}(b-r)^2\}q_{\scriptscriptstyle m}\\
		&=-2\gamma\lambda\phi r(a_{\scriptscriptstyle m}-1)(b-r).\label{eq:MFq2}
	\end{align}

\begin{assumption}\label{assum:mf}
		$\lambda >
		2r+2\rho_{\scriptscriptstyle m} p_{\scriptscriptstyle m}(b-r)+\tfrac{1}{2}\rho_{\scriptscriptstyle m} p_{\scriptscriptstyle m}^2$.
	\end{assumption}

\begin{align}\label{eq:M-Hpifinal}
\widehat{\pi}(x,\bar{x}) =\rho_{\scriptscriptstyle m} p_{\scriptscriptstyle m} x+\frac{\rho_{\scriptscriptstyle m}q_{\scriptscriptstyle m}+\delta_{\scriptscriptstyle m}\rho_{\scriptscriptstyle m} p_{\scriptscriptstyle m}}{1-\delta_{\scriptscriptstyle m}}\bar{x}+\frac{k_{3,m}}{1-\delta_{\scriptscriptstyle m}},
\end{align}

\begin{theorem}
Let \(Q_{\scriptscriptstyle m}\neq 0\), and the average processes \(\overline{\underline{br}\widehat{\pi}}(t)\) and  \(\overline{\sigma\widehat{\pi}}(t)\) be given. Then, under Assumption \ref{assum:mf}, if the following candidate investment strategy ${\bpi^o} = \{{\pi^o}\}$ is admissible, that is, ${\bpi^o} \in \boldsymbol{\Pi}^m$,
\begin{align}\label{eq:M-pi_i}
{\pi^o}(x, \bar{x}) = \rho_{\scriptscriptstyle m} p_{\scriptscriptstyle m} x + \rho_{\scriptscriptstyle m} q_{\scriptscriptstyle m} \bar{x} + k_{1,m} \overline{\sigma\widehat{\pi}} + k_{2,m}\overline{\underline{br}\widehat{\pi}}+k_{3,m},
\end{align}
then, the representative agent's equilibrium investment strategy \(\widehat{\bpi}\) equals ${\bpi^o}$,
in which \(\rho_{\scriptscriptstyle m}\) is given in \eqref{eq:MFrho},  \(p_{\scriptscriptstyle m}=z_{\scriptscriptstyle m}/[\rho_{\scriptscriptstyle m}(b-r)]\) with \(z_{\scriptscriptstyle m}\) solving  \eqref{eq:MF-p1}, \(q_{\scriptscriptstyle m}\) is the solution of \eqref{eq:MFq2}, and
			\(k_j\) $($for \(j=1,2,3\)$)$ are given in \eqref{eq:MFk1}, \eqref{eq:MFk2}, and \eqref{eq:MFk3}, respectively.
			Moreover, the corresponding equilibrium value function is given by
			\begin{align}\label{eq:MF-VV}
				V(x,\bar{x})=A_{\scriptscriptstyle m} x^2+C_{\scriptscriptstyle m} \bar{x}^2+D_{\scriptscriptstyle m} x\bar{x}+E_{\scriptscriptstyle m} x+F_{\scriptscriptstyle m} \bar{x}+I_{\scriptscriptstyle m},
			\end{align}
			in which
			\begin{align}
				A_{\scriptscriptstyle m}&=\mu_1 a_{\scriptscriptstyle m}+\gamma {a^2_{\scriptscriptstyle m}}{-\gamma  \tilde{a}_{\scriptscriptstyle m},}\\
				C_{\scriptscriptstyle m}&=\frac{\lambda\gamma(c_{\scriptscriptstyle m}+\phi)-\frac{\rho_{\scriptscriptstyle m}}{2}\gamma \tilde{a}_{\scriptscriptstyle m}{q^2_{\scriptscriptstyle m}}}{2r-\lambda},\\
				E_{\scriptscriptstyle m}&=\frac{\mu_1\alpha_{\scriptscriptstyle m}r-D_{\scriptscriptstyle m}\overline{\underline{br}\widehat{\pi}}-\rho_{\scriptscriptstyle m}\gamma  \tilde{a}_{\scriptscriptstyle m}\varrho_{\scriptscriptstyle m}p_{\scriptscriptstyle m}+2(a_{\scriptscriptstyle m}-1)\gamma\lambda \alpha_{\scriptscriptstyle m}-\mu_2\lambda}{r-\lambda},\\
				F_{\scriptscriptstyle m}&=\frac{2\lambda\gamma \alpha_{\scriptscriptstyle m}(c_{\scriptscriptstyle m}+\phi)+\lambda\mu_2\phi-\rho_{\scriptscriptstyle m}\gamma \tilde{a}_{\scriptscriptstyle m}\varrho_{\scriptscriptstyle m}q_{\scriptscriptstyle m}-2C_{\scriptscriptstyle m}\overline{\underline{br}\widehat{\pi}}}{r-\lambda},\\
				I_{\scriptscriptstyle m}&=\frac{(C_{\scriptscriptstyle m}-\gamma {c^2_{\scriptscriptstyle m}})\overline{\sigma\widehat{\pi}}^2-\lambda \gamma {\alpha^2_{\scriptscriptstyle m}}+{\frac{\rho_{\scriptscriptstyle m}}{2}\gamma \tilde{a}_{\scriptscriptstyle m}\varrho^2_{\scriptscriptstyle m}}+F_{\scriptscriptstyle m}\overline{\underline{br}\widehat{\pi}}}{\lambda},
			\end{align}
			with
			\begin{align}
				\alpha_{\scriptscriptstyle m}&=\la^{-1}\big({c_{\scriptscriptstyle m}\overline{\underline{br}\widehat{\pi}}+a_{\scriptscriptstyle m}\rho_{\scriptscriptstyle m} \varrho_{\scriptscriptstyle m}(b-r)}\big), \label{eq:MFe1}\\
				\varrho_{\scriptscriptstyle m}&={\left[k_1\overline{\sigma\widehat{\pi}}+k_2\overline{\underline{br}\widehat{\pi}}+k_3\right]/\rho_{\scriptscriptstyle m}}.
			\end{align}
	\end{theorem}

	\begin{proof}
We postulate that for any fixed realization \(\omega \in \Omega\), the average processes \(\overline{\sigma\widehat{\pi}}\) and \(\overline{\underline{br}\widehat{\pi}}\) take the following affine form:
	\begin{align}\label{eq:average_pi}
		\overline{\sigma\hat{\pi}}  = \mathfrak{H}_1\bar{x}  + \mathfrak{H}_0,
		\qquad
		\overline{\underline{br}\hat{\pi}} = \mathfrak{T}_1\bar{x}  + \mathfrak{T}_0,
	\end{align}
	where  $\mathfrak{H}_i$ and \(\mathfrak{T}_i\) (for \(i=0,1\)) are constants and need to be determined.
Then, by substituting \eqref{eq:average_pi} into \eqref{eq:M-pi_i}, we can rewrite the expression of \(\widehat{\pi}(x,\bar{x})\) as follows
\begin{align}
\widehat{\pi}(x, \bar{x})&= \rho_{\scriptscriptstyle m} p_{\scriptscriptstyle m} x + \rho_{\scriptscriptstyle m} q_{\scriptscriptstyle m} \bar{x} + k_{1,\scriptscriptstyle m} \left(\mathfrak{H}_1\,\bar{x} +\mathfrak{H}_0\right)+ k_{2,\scriptscriptstyle m}\left(\mathfrak{T}_1\,\bar{x} + \mathfrak{T}_0\right) +k_{3,\scriptscriptstyle m}\\
		&=\rho_{\scriptscriptstyle m}p_{\scriptscriptstyle m} x+\left(\rho_{\scriptscriptstyle m} q_{\scriptscriptstyle m} +k_{1,\scriptscriptstyle m}\mathfrak{H}_1+k_{2,\scriptscriptstyle m}\mathfrak{T}_1\right)\bar{x}+k_{1,\scriptscriptstyle m}\mathfrak{H}_0+k_{2,\scriptscriptstyle m}\mathfrak{T}_0+k_{3,\scriptscriptstyle m}. \label{eq:M-pi}
	\end{align}
By multiplying $\sigma$ or $(b-r)$ on both sides of \eqref{eq:M-pi}, and taking the conditional expectation, we obtain
\begin{align*}
\overline{\sigma\widehat{\pi}}&=\sigma\left(\rho_{\scriptscriptstyle m}p_{\scriptscriptstyle m}+\rho_{\scriptscriptstyle m} q_{\scriptscriptstyle m} +k_{1,\scriptscriptstyle m}\mathfrak{H}_1+k_{2,\scriptscriptstyle m}\mathfrak{T}_1\right)\bar{x}+\sigma \left(k_{1,\scriptscriptstyle m}\mathfrak{H}_0+k_{2,\scriptscriptstyle m}\mathfrak{T}_0+k_{3,\scriptscriptstyle m}\right),\\
\overline{\underline{br}\widehat{\pi}}&=(b-r)\left(\rho_{\scriptscriptstyle m}p_{\scriptscriptstyle m}+\rho_{\scriptscriptstyle m} q_{\scriptscriptstyle m} +k_{1,\scriptscriptstyle m}\mathfrak{H}_1+k_{2,\scriptscriptstyle m}\mathfrak{T}_1\right)\bar{x}+(b-r) \left(k_{1,\scriptscriptstyle m}\mathfrak{H}_0+k_{2,\scriptscriptstyle m}\mathfrak{T}_0+k_{3,\scriptscriptstyle m}\right).
\end{align*}
By comparing them with \eqref{eq:average_pi}, we deduce the following system of equations
\begin{align}
\begin{cases}
\mathfrak{H}_1&=\sigma\left(\rho_{\scriptscriptstyle m}p_{\scriptscriptstyle m}+\rho_{\scriptscriptstyle m} q_{\scriptscriptstyle m} +k_{1,\scriptscriptstyle m}\mathfrak{H}_1+k_{2,\scriptscriptstyle m}\mathfrak{T}_1\right),\\
\mathfrak{H}_0&=\sigma \left(k_{1,\scriptscriptstyle m}\mathfrak{H}_0+k_{2,\scriptscriptstyle m}\mathfrak{T}_0+k_{3,\scriptscriptstyle m}\right),\\
\mathfrak{T}_1&=(b-r)\left(\rho_{\scriptscriptstyle m}p_{\scriptscriptstyle m}+\rho_{\scriptscriptstyle m} q_{\scriptscriptstyle m} +k_{1,m}\mathfrak{H}_1+k_{2,m}\mathfrak{T}_1\right),\\
\mathfrak{T}_0&=(b-r) \left(k_{1,\scriptscriptstyle m}\mathfrak{H}_0+k_{2,\scriptscriptstyle m}\mathfrak{T}_0+k_{3,m}\right).
\end{cases}
\end{align}
Given that $1-\delta_{\scriptscriptstyle m}\neq 0$, we solve the above system of equations and obtain
	\begin{equation}
		\begin{aligned}
			\mathfrak{H}_1&=\dfrac{\rho_{\scriptscriptstyle m} \sigma (p_{\scriptscriptstyle m}+q_{\scriptscriptstyle m})}{1 - \delta_m},\qquad
			\mathfrak{T}_1=\dfrac{\rho_{\scriptscriptstyle m}(b-r)(p_{\scriptscriptstyle m}+q_{\scriptscriptstyle m})}{1-\delta_m (b-r)},\\
			\mathfrak{H}_0&= \dfrac{\sigma k_{3,\scriptscriptstyle m}}{1-\delta_m}, \qquad \qquad \qquad
			\mathfrak{T}_0=\dfrac{(b-r) k_{3,\scriptscriptstyle m}}{1-\delta_m}.
		\end{aligned}
	\end{equation}
	By substituting these coefficients into \eqref{eq:M-pi} and rearranging terms, we derive the explicit expressions of $\widehat{\pi}$ given in \eqref{eq:M-Hpifinal}.
To show the self-consistency of the obtained ansatz. Define
	\begin{align*}
		\mathfrak{G}_{\scriptscriptstyle m}&:=\frac{\rho_{\scriptscriptstyle m}q_{\scriptscriptstyle m}+\delta_{\scriptscriptstyle m}\rho_{\scriptscriptstyle m} p_{\scriptscriptstyle m}}{1-\delta_{\scriptscriptstyle m}},\\
		\mathfrak{R}_{\scriptscriptstyle m}&:=\frac{k_{3,m}}{1-\delta_{\scriptscriptstyle m}}.
	\end{align*}
	Multiplying \eqref{eq:M-Hpifinal} by $\sigma$ or $(b-r)$ and taking the conditional expectation on both sides of the equation, it yields that
	\begin{align*}
		&\sigma \rho_{\scriptscriptstyle m} p_{\scriptscriptstyle m}+  \sigma\mathfrak{G}_{\scriptscriptstyle m}=\mathfrak{H}_1,\qquad \qquad \qquad\sigma \mathfrak{R}_{\scriptscriptstyle m}=\mathfrak{H}_0,\\
&(b-r) \rho_{\scriptscriptstyle m} p_{\scriptscriptstyle m}+  (b-r)\mathfrak{G}_{\scriptscriptstyle m}=\mathfrak{T}_1,\quad (b-r) \mathfrak{R}_{\scriptscriptstyle m}=\mathfrak{T}_0.
	\end{align*}
Hence, we complete our proof.
\end{proof}

\end{document}
